\bindcontent{
\section*{Study Topic 1: Foundations of Statistics}

\begin{concept}{Basic Statistical Terminology}
    Statistics is the science of collecting, describing, and interpreting data. A \textbf{population} is the entire group of interest, and a \textbf{parameter} describes it. A \textbf{sample} is a subset of the population, and a \textbf{statistic} describes the sample.
\end{concept}

\begin{concept}{Types of Data}
    \textbf{Qualitative (categorical)} variables classify individuals based on attributes. \textbf{Quantitative} variables provide numerical measures and can be \textbf{discrete} (countable) or \textbf{continuous} (measurable).
\end{concept}

\begin{concept}{Levels of Measurement}
    \begin{enumerate}[label=\roman*.]
        \item \textbf{Nominal:} Names or labels.
        \item \textbf{Ordinal:} Ordered categories.
        \item \textbf{Interval:} Meaningful differences, no true zero.
        \item \textbf{Ratio:} Meaningful ratios, true zero exists.
    \end{enumerate}
\end{concept}

\begin{concept}{Study Design}
    \textbf{Observational studies} measure characteristics without influence, showing association. \textbf{Designed experiments} apply treatments to observe effects, allowing for causation claims.
\end{concept}

\begin{concept}{Experimental Design}
    Observational studies cannot prove causation due to \textbf{lurking or confounding variables} that may affect the response variable.
\end{concept}

\begin{concept}{Sampling Methods}
    In a \textbf{simple random sample}, every individual and every possible sample of size $n$ has an equal chance of being selected. Use \texttt{randInt(min, max, n)} on a calculator to generate.
\end{concept}
}
